\documentclass[11pt,a4paper]{article}
\usepackage[utf8]{inputenc}
\usepackage[T1]{fontenc}
\usepackage{lmodern}
\usepackage{textcomp}
\usepackage{amsmath,amssymb}
\usepackage{graphicx}
\usepackage{booktabs}
\usepackage{longtable}
\usepackage{array}
\usepackage{listings}
\usepackage{xcolor}
\usepackage[margin=1in]{geometry}
\usepackage{float}
\usepackage[hidelinks,breaklinks]{hyperref}
\usepackage{microtype}
\setlength{\parskip}{4pt}
\lstset{basicstyle=\ttfamily\scriptsize,breaklines=true,breakatwhitespace=false,
  postbreak=\mbox{\textcolor{gray}{$\hookrightarrow$}\space},frame=single,
  framerule=0.3pt,rulecolor=\color{gray!60},columns=fullflexible,keepspaces=true,
  showstringspaces=false,numbers=none,xleftmargin=4pt,xrightmargin=4pt,
  literate={\_}{{\_}}1}
\graphicspath{{figures/}}
\providecommand{\tightlist}{}
\title{Transcript: The Thermal Pure Higgsino near 1.1 TeV versus Planck 2018, LZ 2024 and XENONnT: A Literature-Based Assessment}
\author{HEP-Genesis-Agent session record \\ \small Project: ramp-research-higgsino-dd}
\date{2026-10-16}
\begin{document}
\maketitle
\noindent\textit{\small The conversation verbatim, generated mechanically from the session record.}
\vspace{1em}



\appendix
\section{Conversation transcript}
\label{app:transcript}
The exchange verbatim (assistant text as shown to the user; tool calls summarized by name --- see the tool-call log for arguments and results).
\subsection*{User \small(2026-10-16 11:03:59)}
Survey recent arXiv work on the higgsino thermal relic target (1.0-1.2 TeV). Use the inelastic-DM tools to compute the relic density at several masses; compare with LZ/XENONnT sensitivities and the Planck abundance. Where does the higgsino sit relative to current limits?

\subsection*{Assistant \small(research mode, 2026-10-16 11:03:59)}
\subsection*{User \small(2026-10-16 11:04:37)}
Survey recent arXiv work on the higgsino thermal relic target (1.0-1.2 TeV). Use the inelastic-DM tools to compute the relic density at several masses; compare with LZ/XENONnT sensitivities and the Planck abundance. Where does the higgsino sit relative to current limits?

\subsection*{Assistant \small(research mode, 2026-10-16 11:04:37)}
\textit{Plan:}\begin{enumerate}
  \item (closed) Establish the pure-higgsino thermal relic density as a function of mass (about 0.9, 1.0, 1.1, 1.2 and 1.3 TeV), with and without Sommerfeld enhancement, and find the mass at which Omega h\textasciicircum{}2 equals the Planck 2018 value of 0.120 +/- 0.001.
  \item (closed) Place the pure-higgsino elastic spin-independent cross section at about 1.0-1.2 TeV against the LZ 2024 (2410.17036) and XENONnT (2303.14729) limits at the same mass, and against the xenon neutrino fog.
  \item (closed) Assess the routes by which a higgsino becomes visible to direct detection: the mass splittings and the inelastic Z-exchange threshold, LZ inelastic constraints (2609.09830), mixing-induced Higgs couplings in the (mu, M1/M2, tan beta) plane (2606.31013), and the (Omega h\textasciicircum{}2/0.12) rescaling for sub-thermal higgsinos.
  \item (blocked) Pin down the pure-higgsino thermal mass from at least two independent published calculations, including recent work with bound-state formation, and quote the Planck 2018 Omega\_c h\textasciicircum{}2 verbatim. Confirm which mass the 1411.6930 pure-higgsino values (0.126/0.127) refer to.
  \item (blocked) Replace the linear 40 GeV extrapolation with direct reads of the LZ 2024 and XENONnT SI limits at about 1 TeV, a numeric xenon neutrino-fog value at about 1 TeV, and one next-generation projection (XLZD/DARWIN). Correct the f-008 arithmetic and quote Chen and Hill verbatim.
  \item (closed) Verify 2606.31013 and 2609.09830 with metadata, and trace the contested f-002 and f-003 numbers to quoted excerpts: Eq. 13 (356 MeV), the f\_q formula, Eq. 38, the Table VII NLO rows, and the tree-level delta\_m0 formula (Eq. 9 / Ref. [28]).
  \item (open) Identify Ref. [28] of arXiv:2609.09830 and quote its Eq. (9) for the neutral splitting
  \item (open) Add a sourced statement on indirect-detection limits for TeV higgsinos
\end{enumerate}
\textit{Tool calls:} \texttt{list\_agent\_files}, \texttt{list\_input\_files}, \texttt{list\_datasets}, \texttt{list\_files}, \texttt{get\_dispatch}, \texttt{search\_arxiv\_abstracts}, \texttt{search\_arxiv\_abstracts}, \texttt{search\_arxiv\_abstracts}, \texttt{search\_arxiv\_abstracts}, \texttt{search\_arxiv}, \texttt{search\_arxiv}, \texttt{search\_arxiv}, \texttt{search\_arxiv}, \texttt{get\_arxiv\_metadata}, \texttt{search\_arxiv\_abstracts}, \texttt{search\_arxiv\_abstracts}, \texttt{search\_arxiv\_abstracts}, \texttt{list\_input\_files}, \texttt{download\_arxiv\_full\_text}, \texttt{download\_arxiv\_full\_text}, \texttt{search\_arxiv\_abstracts}, \texttt{download\_arxiv\_full\_text}, \texttt{search\_arxiv\_abstracts}, \texttt{download\_arxiv\_full\_text}, \texttt{download\_arxiv\_full\_text}, \texttt{read\_arxiv\_chunk}, \texttt{read\_arxiv\_chunk}, \texttt{read\_arxiv\_chunk}, \texttt{read\_arxiv\_chunk}, \texttt{read\_arxiv\_chunk}, \texttt{read\_arxiv\_chunk}, \texttt{read\_arxiv\_chunk}, \texttt{read\_arxiv\_chunk}, \texttt{read\_arxiv\_chunk}, \texttt{read\_arxiv\_chunk}, \texttt{read\_arxiv\_chunk}, \texttt{read\_arxiv\_chunk}, \texttt{search\_arxiv\_abstracts}, \texttt{read\_arxiv\_chunk}, \texttt{get\_paper\_figures}, \texttt{search\_arxiv\_abstracts}, \texttt{read\_arxiv\_chunk}, \texttt{search\_arxiv\_abstracts}, \texttt{get\_arxiv\_metadata}, \texttt{read\_arxiv\_chunk}, \texttt{download\_arxiv\_full\_text}, \texttt{download\_arxiv\_full\_text}, \texttt{post\_finding}, \texttt{post\_finding}, \texttt{post\_finding}, \texttt{download\_arxiv\_full\_text}, \texttt{read\_arxiv\_chunk}, \texttt{cite}, \texttt{cite}, \texttt{post\_question}, \texttt{post\_finding}, \texttt{post\_finding}, \texttt{read\_arxiv\_chunk}, \texttt{read\_arxiv\_chunk}, \texttt{post\_finding}, \texttt{post\_finding}, \texttt{read\_arxiv\_chunk}, \texttt{post\_finding}, \texttt{cite}, \texttt{cite}, \texttt{cite}, \texttt{post\_question}, \texttt{read\_arxiv\_chunk}, \texttt{post\_finding}, \texttt{get\_arxiv\_metadata}, \texttt{get\_arxiv\_metadata}, \texttt{read\_arxiv\_chunk}, \texttt{get\_paper\_figures}, \texttt{search\_arxiv\_abstracts}, \texttt{get\_paper\_figures}, \texttt{read\_arxiv\_chunk}, \texttt{read\_arxiv\_chunk}, \texttt{read\_arxiv\_chunk}, \texttt{read\_arxiv\_chunk}, \texttt{read\_arxiv\_chunk}, \texttt{read\_arxiv\_chunk}, \texttt{read\_arxiv\_chunk}, \texttt{search\_arxiv\_abstracts}, \texttt{read\_arxiv\_chunk}, \texttt{read\_arxiv\_chunk}, \texttt{search\_arxiv\_abstracts}, \texttt{get\_arxiv\_metadata}, \texttt{read\_arxiv\_chunk}, \texttt{get\_arxiv\_metadata}, \texttt{read\_arxiv\_chunk}, \texttt{post\_finding}, \texttt{post\_finding}, \texttt{post\_dead\_end}, \texttt{post\_question}, \texttt{post\_finding}, \texttt{post\_finding}, \texttt{post\_question}, \texttt{post\_dead\_end}, \texttt{post\_finding}, \texttt{post\_finding}, \texttt{post\_finding}, \texttt{post\_finding}, \texttt{post\_question}\par
\textbf{Survey} --- 29 tools on 2 servers, 5 read-only probes, 11 papers found.

The literature agrees that a pure higgsino (an SU(2) doublet with Y=1/2) makes up all of the dark matter as a thermal relic at a mass of roughly 1.0-1.2 TeV. Sommerfeld-enhanced coannihilation calculations, for example 1411.6930, move the exact mass within that window, and 2606.31013 quotes about 1 TeV. Pure-higgsino scattering in direct detection comes from loops and is tiny: the one- and two-loop electroweak calculations (1104.0228 and 1007.2601) put the spin-independent cross section well below current xenon limits, near or inside the neutrino fog. I did not read those full texts, so the exact numbers need checking. The 2024 LZ result (2410.17036, 4.2 tonne-years) and the first XENONnT result (2303.14729) are tightest near 30-40 GeV. At 1 TeV their limits scale up roughly linearly with mass, to about 1e-46 to 1e-45 cm\textasciicircum{}2 (my estimate, not read from their tables). A thermal pure higgsino is therefore not excluded by elastic scattering. Direct detection can see it only if mixing with gauginos raises the Higgs coupling and the mass splittings (2606.31013), or if the splitting between the two neutral states is small enough (roughly up to a few hundred keV) that tree-level inelastic Z-exchange turns on. That inelastic case is already strongly constrained, and 2609.09830 recomputes it at one loop for LZ. Planck 2018 (1807.06209) sets the target abundance at Omega\_c h\textasciicircum{}2 = 0.120 +/- 0.001.

Quantities worth checking: Omega h\textasciicircum{}2 against higgsino mass at about 0.9, 1.0, 1.1, 1.2 and 1.3 TeV, with and without Sommerfeld enhancement; find the mass where Omega h\textasciicircum{}2 = 0.120; Charged-neutral splitting (about 340 MeV from electroweak loops in the pure limit) and neutral-neutral splitting delta\_m0 as functions of the gaugino masses M1 and M2; Elastic spin-independent cross section of the pure higgsino at 1.1 TeV compared with the LZ 2024 and XENONnT limits at the same mass, and with the xenon neutrino fog; Inelastic Z-exchange threshold: the largest delta\_m0 at which xenon at about 1 TeV can still produce a recoil (scale set by the kinematic limit, about mu v\textasciicircum{}2 / 2), and the delta\_m0 values LZ excludes; Rescaling of signal rate by (Omega h\textasciicircum{}2 / 0.12) for higgsinos lighter than the thermal mass that make up only part of the dark matter; Size of the tree-level Higgs coupling as a function of M1, M2 and tan(beta), to map where mixing makes higgsinos visible to LZ.

\textbf{Plan} --- 3 lines of investigation:

\begin{enumerate}
  \item Establish the pure-higgsino thermal relic density as a function of mass (about 0.9, 1.0, 1.1, 1.2 and 1.3 TeV), with and without Sommerfeld enhancement, and find the mass at which Omega h\textasciicircum{}2 equals the Planck 2018 value of 0.120 +/- 0.001. --- This is the core of the task. The catalog has no inelastic-DM or relic-density compute tool and no coding-task tool, so the Omega h\textasciicircum{}2(m) curve has to come from published calculations. Each value must be read from text, tables or figures, never estimated.
  \item Place the pure-higgsino elastic spin-independent cross section at about 1.0-1.2 TeV against the LZ 2024 (2410.17036) and XENONnT (2303.14729) limits at the same mass, and against the xenon neutrino fog. --- This answers the question of where the higgsino sits relative to current limits. The survey's 1e-46 to 1e-45 cm\textasciicircum{}2 limit at 1 TeV is an unverified estimate and must be replaced with read values.
  \item Assess the routes by which a higgsino becomes visible to direct detection: the mass splittings and the inelastic Z-exchange threshold, LZ inelastic constraints (2609.09830), mixing-induced Higgs couplings in the (mu, M1/M2, tan beta) plane (2606.31013), and the (Omega h\textasciicircum{}2/0.12) rescaling for sub-thermal higgsinos. --- Elastic scattering leaves the pure higgsino invisible, so whether it is visible depends on the splittings set by gaugino masses. This direction goes past the base question and is a robustness check on any 'not excluded' conclusion.
Critic's view: revise --- Step 1 is not enough: the gap has to show up in the final answer. No relic-density or inelastic-DM compute tool exists, list\_input\_files is empty, and run\_test\_job only draws random samples. The task asked us to compute relic density with tools, and we cannot. The report must say plainly that every Omega h\textasciicircum{}2 value is taken from the literature. Do not use run\_test\_job to produce anything that looks like a computed relic curve.; 1411.6930 is a pMSSM scan. It may not give a pure-higgsino Omega h\textasciicircum{}2(m) curve at 0.9-1.3 TeV with and without Sommerfeld enhancement. Before relying on it, check that such a figure or table exists. If it does not, use the one-parameter thermal mass plus the m\textasciicircum{}2 scaling, and label every resulting row as derived, not read. Also note that the higgsino Sommerfeld effect is small, only a few percent shift in mass. Bound-state formation (recent 2021-2026 work) can matter as much as Sommerfeld, so the search should target it explicitly.; The budget is tight. Fifteen calls do not cover three or four full-text downloads plus chunk reads, figures and cite. Drop the full read of 1008.2905, which is a method paper and not needed for numbers. Get the Planck Omega\_c h\textasciicircum{}2 = 0.120 +/- 0.001 quote with one targeted chunk read of 1807.06209. Limit independent cross-checks to one paper..

\end{enumerate}
\textbf{Wave 1} --- now exploring, in parallel: Establish the pure-higgsino thermal relic density as a function of mass (about 0.9, 1.0, 1.1, 1.2 and 1.3 TeV), with and without Sommerfeld enhancement, and find the mass at which Omega h\textasciicircum{}2 equals the Planck 2018 value of 0.120 +/- 0.001.; Place the pure-higgsino elastic spin-independent cross section at about 1.0-1.2 TeV against the LZ 2024 (2410.17036) and XENONnT (2303.14729) limits at the same mass, and against the xenon neutrino fog.; Assess the routes by which a higgsino becomes visible to direct detection: the mass splittings and the inelastic Z-exchange threshold, LZ inelastic constraints (2609.09830), mixing-induced Higgs couplings in the (mu, M1/M2, tan beta) plane (2606.31013), and the (Omega h\textasciicircum{}2/0.12) rescaling for sub-thermal higgsinos..

\textbf{Finding:} Kinematic threshold for inelastic up-scattering on xenon (my arithmetic): m\_A = 131 x 0.9315 GeV = 122.0 GeV; mu\_A = 108.7 GeV (m=1.0 TeV) and 109.8 GeV (m=1.1 TeV); v\_max = 544+232 = 776 km/s, beta\textasciicircum{}2 = (776/299792)\textasciicircum{}2 = 6.70e-6; delta\_max = mu\_A beta\textasciicircum{}2 /2 = 364 keV (1.0 TeV) and 368 keV (1.1 TeV). The minimum-velocity formula v\_min = sqrt(2 delta/mu\_A) is given in 2609.09830 Appendix A (Eq. A7).

\textbf{Finding:} 2609.09830 does not quote an LZ-excluded delta\_m0 range. Its five benchmarks (\textbar{}mu\textbar{} = 0.61-1.40 TeV; \textbar{}M1\textbar{}, M2 about 5.8-42 TeV; tan beta 10.9-26.5) have one-loop on-shell splittings delta = 8.9-36.4 keV. Each predicts at most about 1.03 combined NLO events in LZ, and probing them needs 9.4-65.5 tonne-years. All are consistent with the LZ bound. Inelastic events at NLO range from 0.00022 (BP1) to 0.18 (BP4). NLO vertex corrections can exceed 100 percent, and in BP2 the inelastic cross section is about twice the LO elastic one. Splittings of O(10) keV come from partial cancellation of the M\_Z\textasciicircum{}2/M1 and M\_Z\textasciicircum{}2/M2 terms when M1 and M2 have opposite signs.

\textbf{Finding:} 2606.31013: the SI coupling comes from Higgs exchange, f\_q \textasciitilde{} -(g\textasciicircum{}2/8M\_h\textasciicircum{}2)[1/(M2-\textbar{}mu\textbar{}) + tan\textasciicircum{}2 theta\_W/(M1-\textbar{}mu\textbar{})](1 + sgn(mu) sin 2beta), so it is suppressed at large gaugino mass and at small tan beta with mu\textless{}0. The scan covers 150\textless{}\textbar{}mu\textbar{}\textless{}1000 GeV, 1.6\textless{}tan beta\textless{}50, 0.5\textless{}M1\textless{}10 TeV with M2=2M1. Rates are rescaled by Omega\_H h\textasciicircum{}2/0.12 (Eq. 38, computed with micrOMEGAs 6.3.0). Direct detection is sensitive mainly where higgsino-gaugino mixing produces sizable charged-neutral splittings, and reaches the thermal mass of about 1 TeV. The radiative charged-neutral splitting tends to about 356 MeV for M \textgreater{}\textgreater{} M\_Z (Eq. 13).

\textbf{Finding:} Hisano et al. (1104.0228) find the leading-order electroweak-loop SI cross section of pure EW-multiplet DM (including the higgsino, n=2, Y=1/2, evaluated at M=1 TeV) to be O(1e-46 to 1e-48) cm\textasciicircum{}2, suppressed by an accidental cancellation: twist-2 contributions are positive while Higgs-exchange, gluon and scalar contributions are negative, and they are of similar size. The cross section is nearly independent of mass above about 1 TeV.

\textbf{Finding:} LZ 2024 (2410.17036, 4.2 tonne-years): the strongest SI exclusion is 2.2e-48 cm\textasciicircum{}2 (90\% CL) and the best median sensitivity is 5.1e-48 cm\textasciicircum{}2, both at 40 GeV/c\textasciicircum{}2. The abstract quotes no value at 1 TeV.

\textbf{Finding:} No relic-density or inelastic-DM compute tool and no input data tables are available: list\_input\_files returned an empty file list, and the catalog contains only statistics/histogram, test-job and literature tools. All Omega h\textasciicircum{}2 values for d-01 therefore come from the literature and were not computed.

\textbf{Finding:} Beneke, Hellmann and Ruiz-Femenia (arXiv:1411.6930, Sec. 3) report Omega h\textasciicircum{}2 = 0.127 (perturbative) and 0.126 (Sommerfeld, with off-diagonal rates) for a pure-higgsino model compared against the higgsino-like pMSSM point 1627006 at m\_chi = 1172.31 GeV. The net Sommerfeld correction for pure higgsinos is almost zero (below 1 percent) because enhancement in the neutral channels cancels suppression in the chi0 chi+- channels. For the pMSSM higgsino-like point (higgsino fraction 0.98, dm+ = 1.8 GeV, dm0 = 9.5 GeV) they find Omega h\textasciicircum{}2 = 0.108 (perturbative) and 0.100 (Sommerfeld), an 8 percent reduction. The quoted Sommerfeld accuracy is about 5 percent at best.

\textbf{Completed --- Assess the routes by which a higgsino becomes visible to direct detection: the mass splittings and the inelastic Z-exchange threshold, LZ inelastic constraints (2609.09830), mixing-induced Higgs couplings in the (mu, M1/M2, tan beta) plane (2606.31013), and the (Omega h\textasciicircum{}2/0.12) rescaling for sub-thermal higgsinos..} Two of the three success criteria are met; the third is only partly met. (1) The inelastic threshold, worked by hand (f-001): m\_Xe = 131 x 0.9315 = 122.0 GeV. The reduced mass mu\_A is 108.7 GeV at m = 1.0 TeV and 109.8 GeV at 1.1 TeV. With v\_max = 544 + 232 = 776 km/s, (v\_max/c)\textasciicircum{}2 = 6.70e-6, so delta\_max = mu\_A v\_max\textasciicircum{}2/2 is about 364 keV at 1.0 TeV and 368 keV at 1.1 TeV. The minimum-velocity formula v\_min = sqrt(2 delta/mu\_A) is Eq. A7 of 2609.09830. (2) Partly met: 2609.09830 does not give an LZ-excluded delta\_m0 range (f-002). Its five benchmarks have \textbar{}mu\textbar{} = 0.61-1.40 TeV, \textbar{}M1\textbar{} and M2 of about 5.8-42 TeV, and one-loop on-shell delta = 8.9-36.4 keV. They predict at most about 1.03 combined NLO events in LZ, need 9.4-65.5 t-y of exposure to be probed (Table VII, Fig. 6), and are all consistent with LZ. NLO vertex corrections can exceed 100 percent, and in BP2 the inelastic cross section is about twice the LO elastic one. Splittings of O(10) keV with gaugino masses of O(10) TeV arise only from partial cancellation between the M\_Z\textasciicircum{}2/M1 and M\_Z\textasciicircum{}2/M2 terms when M1 and M2 have opposite signs; generically the paper quotes O(10-100) MeV. My own estimate, using the formula that is not yet sourced (q-05): same-sign M2 = 2 M1 would need M1 of about 1.4e7 GeV to bring delta below 364 keV. Inelastic Z-exchange therefore matters only in tuned corners. (3) Mixing region where elastic SI becomes visible (f-003, 2606.31013): Higgs exchange gives f\_q \textasciitilde{} -(g\textasciicircum{}2/8M\_h\textasciicircum{}2)[1/(M2-\textbar{}mu\textbar{}) + tan\textasciicircum{}2 theta\_W/(M1-\textbar{}mu\textbar{})](1 + sgn(mu) sin 2beta). This is suppressed at large gaugino mass and at small tan beta with mu\textless{}0. The scan covers 150\textless{}\textbar{}mu\textbar{}\textless{}1000 GeV, 1.6\textless{}tan beta\textless{}50 and 0.5\textless{}M1\textless{}10 TeV with M2 = 2 M1. Direct detection is sensitive where mixing produces larger charged-neutral splittings and reaches the thermal mass of about 1 TeV (Figs. 4 and 5 of that paper; the figures were not extracted). The pure-limit radiative charged-neutral splitting is about 356 MeV (Eq. 13 of 2606.31013), rather than the 340 MeV in the brief. (4) Rescaling: 2606.31013 Eq. 38 uses sigma\_eff = (Omega\_H h\textasciicircum{}2/0.12) sigma, computed with micrOMEGAs 6.3.0, and 2609.09830 states that limits weaken proportionally for sub-thermal fractions. Quotes from both papers are cited in /Users/nramachandra/hep-genesis-testing/projects/ramp-research-higgsino-dd/runs/261016\_110437/references.bib. Full texts are at the same run directory as 2609.09830.txt and 2606.31013.txt.

\textbf{Finding:} Chen and Hill (1912.07795) include the O(mW/M) power corrections, perturbative QCD to NNNLO and ChEFT nuclear effects. They find that the pure-higgsino SI cross section has an upper limit below 1e-48 cm\textasciicircum{}2 in the TeV mass range, below the estimated xenon neutrino floor. This is about 2 orders of magnitude below the LZ 2024 minimum limit of 2.2e-48 cm\textasciicircum{}2 at 40 GeV. Assuming linear mass scaling of the LZ limit above about 100 GeV (my estimate, not read from the paper), the LZ limit at 1.1 TeV is about 2.2e-48 x (1100/40), roughly 6e-47 cm\textasciicircum{}2, so the pure higgsino lies at least about 1.8 orders of magnitude below it.

\textbf{Completed --- Place the pure-higgsino elastic spin-independent cross section at about 1.0-1.2 TeV against the LZ 2024 (2410.17036) and XENONnT (2303.14729) limits at the same mass, and against the xenon neutrino fog..} Answer: a thermal pure higgsino at about 1.1 TeV lies about 2 or more orders of magnitude below the LZ 2024 limit and below the xenon neutrino floor. Three of the planned inputs were not read: the XENONnT limit, the LZ limit read from its figure at 1 TeV, and a numeric neutrino-fog value. All are recorded as open question q-06.

(1) Prediction. Hisano et al. 1104.0228 work at leading order in electroweak loops. They find the SI cross section of pure SU(2)-multiplet dark matter is O(1e-46 to 1e-48) cm\textasciicircum{}2, nearly independent of mass above about 1 TeV. An accidental cancellation drives this: twist-2 terms are positive, while the Higgs-exchange, gluon and Z-box scalar terms are negative and of similar size. Their Fig. 5 shows the higgsino (n=2, Y=1/2) at M = 1 TeV against the Higgs mass; the paper predates the measured Higgs mass. Finding f-004.

The newer and more complete calculation is Chen and Hill 1912.07795, which adds O(mW/M) power corrections, QCD corrections to NNNLO and ChEFT nuclear effects. Its verbatim result is: 'The upper limit for the pure doublet is less than 10\textasciicircum{}-48 cm\textasciicircum{}2 in the TeV mass regime ... lower than the estimated Xenon neutrino floor', with 'the pure Higgsino well below the neutrino floor'. The uncertainty is quoted only as this upper bound; no central value plus or minus an error was read. Citation recorded in references.bib (tc-...-64).

(2) LZ 2024 limit (2410.17036). Verbatim: the strongest SI exclusion is 2.2e-48 cm\textasciicircum{}2 at 90\% CL, and the best median sensitivity is 5.1e-48 cm\textasciicircum{}2, both at 40 GeV. Finding f-005, citation tc-...-65. No 1 TeV value was read from the text or a figure.

(3) Rescaling. Scaling the 40 GeV limit linearly with mass is my estimate, not a read value. It gives about 2.2e-48 x (1100/40), roughly 6e-47 cm\textasciicircum{}2, at 1.1 TeV. The ratio of limit to predicted upper bound is then at least about 60, i.e. at least 1.8 orders of magnitude. This extrapolation is likely conservative because the limit curve rises from its minimum. Finding f-008, medium confidence.

Not done:

\begin{itemize}
  \item The XENONnT 2303.14729 1 TeV limit was not read.
  \item No limit figure was extracted, so there are no figure numbers for the LZ or XENONnT limits.
  \item No numeric neutrino-fog value at 1 TeV was obtained; the 'below the neutrino floor' statement rests on Chen and Hill's text.
  \item Projected sensitivities (LZ full exposure, XLZD/DARWIN) were not tabulated.
\end{itemize}
Also noted: 2607.20588 (Bisal et al., 2026) says next-to-leading-order corrections can push the higgsino cross section below the neutrino floor or make it vanish in parts of MSSM parameter space. This is relevant to q-04.

Files are in /Users/nramachandra/hep-genesis-testing/projects/ramp-research-higgsino-dd/runs/261016\_110437/: 1104.0228.txt, 2410.17036.txt, 1912.07795.txt and references.bib.

\textbf{Finding (preliminary):} Hisano et al. (hep-ph/0610249) expect the Sommerfeld effect on the higgsino-like relic abundance to be roughly 10 percent, compared with a 50 percent reduction for the wino. 1411.6930 later found the net pure-higgsino effect is below 1 percent (Omega h\textasciicircum{}2 = 0.127 perturbative, 0.126 Sommerfeld). Rescaling that single point (taken as m = 1172.31 GeV) with Omega h\textasciicircum{}2 proportional to m\textasciicircum{}2 gives the Planck value 0.120 at m of about 1.14 TeV (Sommerfeld: 1144 GeV; perturbative: 1139 GeV). The rescaled values are 0.074, 0.092, 0.111, 0.132 and 0.155 at m = 0.9, 1.0, 1.1, 1.2 and 1.3 TeV. These are m\textasciicircum{}2 rescalings of one published point, not published values.

\textbf{Blocked --- Establish the pure-higgsino thermal relic density as a function of mass (about 0.9, 1.0, 1.1, 1.2 and 1.3 TeV), with and without Sommerfeld enhancement, and find the mass at which Omega h\textasciicircum{}2 equals the Planck 2018 value of 0.120 +/- 0.001..} (Tool-call budget used up. No tool can compute relic densities, and 1411.6930 gives only one pure-higgsino mass point, so Omega h\textasciicircum{}2 at several masses and a second independent thermal mass were not obtained. Next steps: read 1807.06209.txt (already downloaded) and cite the 0.120 +/- 0.001 value, then find a paper with Omega h\textasciicircum{}2 for higgsinos at several masses, for example later work that includes bound-state formation.) Not finished: only one published pure-higgsino relic-density point was found, so the success criteria (published values at three or more masses, and a thermal mass from two independent papers) are not met.

(1) No relic-density or inelastic-DM compute tool exists here, and list\_input\_files returned no data files (tc-261016\_1-w1-d-01-19, finding f-006). Every value below comes from the literature. Nothing was computed.

(2) arXiv:1411.6930, Sec. 3, gives the one published point: a pure higgsino (all sfermions and heavy Higgses decoupled) has Omega h\textasciicircum{}2 = 0.127 without and 0.126 with Sommerfeld enhancement (off-diagonal rates included). The net Sommerfeld effect is below 1 percent because the neutral channels are enhanced and the charged ones suppressed. The paper's mass for this model appears to be 1172.31 GeV, the mass of its higgsino-like pMSSM point 1627006; I inferred this from context, the paper does not state it. For that pMSSM point (dm+ = 1.8 GeV, dm0 = 9.5 GeV) the values are 0.108 without and 0.100 with Sommerfeld, an 8 percent drop. The paper quotes about 5 percent accuracy at best and compares to Planck 2013 (0.1187 +/- 0.0017). Its Fig. 7 (registered as Figure 6, figures\_1411.6930/fig07.png) shows only the yield ratio Y/Ypert, not Omega h\textasciicircum{}2 against mass (finding f-007). A verbatim quote is in references.bib (key arxiv14116930).

(3) hep-ph/0610249 expected the Sommerfeld effect for higgsinos to be roughly 10 percent (versus 50 percent for winos). It was not a full calculation (f-009).

(4) Interpolation, not published values: scaling the 1411.6930 point as Omega h\textasciicircum{}2 proportional to m\textasciicircum{}2 gives 0.074, 0.092, 0.111, 0.132 and 0.155 at 0.9, 1.0, 1.1, 1.2 and 1.3 TeV. Omega h\textasciicircum{}2 = 0.120 falls at about 1.14 TeV (1144 GeV with Sommerfeld, 1139 GeV without) (f-009, low confidence).

Still to do: the Planck 2018 text was downloaded but not yet read or quoted; 1008.2905 was not read; no independent recent thermal-mass determination was found.

Files in /Users/nramachandra/hep-genesis-testing/projects/ramp-research-higgsino-dd/runs/261016\_110437/: 1411.6930.txt, hep-ph\_0610249.txt, 1807.06209.txt, figures\_1411.6930/fig07.png, references.bib.

\textbf{Critic, wave 1} --- continue: 5 findings stand up to the evidence, 4 questioned ('2606.31013: the SI coupling comes from Higgs exchange, f\_q \textasciitilde{} -(g\textasciicircum{}2/8M\_h\textasciicircum{}2)[1/(M2-\textbar{}mu\textbar{}) + t': None of the three quoted excerpts contains any claimed content. tc-...-35 covers the oblique-parameter chi-square (Eqs. 26-29), tc-...-30 is the title page, and tc-...-33 covers the W/Y parameters (Eqs. 39-40). Missing are the f\_q formula, the scan ranges, Eq. 13 (356 MeV), Eq. 38 with micrOMEGAs 6.3.0, and the statement that direct detection reaches about 1 TeV only with mixing. These may well be in the paper, but the cited evidence does not show them.; 'Beneke, Hellmann and Ruiz-Femenia (arXiv:1411.6930, Sec. 3) report Omega h\textasciicircum{}2 = 0.127 (pert': The excerpts do not show the Omega h\textasciicircum{}2 values 0.127, 0.126, 0.108 or 0.100, nor dm+ = 1.8 GeV, dm0 = 9.5 GeV or the 5 percent accuracy. The claim states m\_chi = 1172.31 GeV as reported, but the d-01 summary says this mass was inferred from context and is not stated in the paper. The figure is also mislabelled: fig07 is registered as Figure 6.; 'Chen and Hill (1912.07795) include the O(mW/M) power corrections, perturbative QCD to NNNL': There are two problems. First, the Chen and Hill verbatim quote (the upper limit below 1e-48 cm\textasciicircum{}2, below the neutrino floor) is not in any quoted excerpt. The cited tc-...-59 and -58 show only the rate formula and the title, and the citation call tc-...-64 is not listed as evidence. Second, there is an arithmetic error. An upper bound below 1e-48 cm\textasciicircum{}2 compared with the 2.2e-48 cm\textasciicircum{}2 minimum is a factor of about 2 or more, not 'about 2 orders of magnitude'. Only the comparison at 1.1 TeV (about 6e-47 cm\textasciicircum{}2) supports a margin of 1.8 orders or more.; '2609.09830 does not quote an LZ-excluded delta\_m0 range. Its five benchmarks (\textbar{}mu\textbar{} = 0.61-': The visible part of Table VII supports only BP1 elastic NLO = 1.028 events, which is consistent with 'at most about 1.03'. The other quantities are not in the shown excerpts: the splitting range 8.9-36.4 keV, the exposures of 9.4-65.5 t-y, the inelastic NLO range, and the BP2 statement. The visible LO numbers for BP2 (inelastic 0.0247, elastic 0.0436) show inelastic below elastic. The 'twice the LO elastic' claim therefore needs the NLO inelastic value quoted explicitly.).

Issues: The core quantity is still unmet. There is no published Omega h\textasciicircum{}2(m) curve at multiple masses and no second, independent determination of the thermal mass. The m\textasciicircum{}2 rescaling in f-009 depends on an inferred mass of 1172.31 GeV. If that inference is wrong, every rescaled value and the 1.14 TeV thermal mass shift with it.; The Planck 2018 value of 0.120 +/- 0.001 is used throughout but has never been quoted from 1807.06209.txt, although the file is already downloaded. 1411.6930 normalised to Planck 2013 (0.1187), so the f-009 rescaling should state which target it uses.; The arithmetic checks out: 0.126 x (m/1172.31)\textasciicircum{}2 gives 0.074, 0.092, 0.111, 0.132 and 0.155, and the thermal masses of 1144 and 1139 GeV are correct. However, the hep-ph/0610249 '10 percent' statement is not visible in the tc-...-68 excerpt..

\textbf{Next steps after wave 1} --- This is the final wave, with 19 tool calls across at most 3 directions. d-01 is closed and replaced by a narrower child that works mostly from files already downloaded: quote Planck 2018, confirm the 1411.6930 mass, and add one independent thermal-mass source that includes bound-state formation. The limits direction carries the headline answer, so it gets priority 1. It reads LZ and XENONnT at 1 TeV from their figures, cites Chen and Hill verbatim, adds fog and XLZD numbers, and fixes the f-008 arithmetic. The third direction is a lower-priority verification pass on the contested f-002 and f-003 evidence and the unsourced delta\_m0 formula (q-02, q-05). No compute tool for relic density exists (f-006, q-01), so the final report must state that all Omega h\textasciicircum{}2 values come from the literature or are labelled derived rescalings. The task does not ask for facility runs, so none is scheduled. Each direction should use about 6 calls and reuse already-downloaded texts wherever possible.

New lines of investigation: Pin down the pure-higgsino thermal mass from at least two independent published calculations, including recent work with bound-state formation, and quote the Planck 2018 Omega\_c h\textasciicircum{}2 verbatim. Confirm which mass the 1411.6930 pure-higgsino values (0.126/0.127) refer to.; Replace the linear 40 GeV extrapolation with direct reads of the LZ 2024 and XENONnT SI limits at about 1 TeV, a numeric xenon neutrino-fog value at about 1 TeV, and one next-generation projection (XLZD/DARWIN). Correct the f-008 arithmetic and quote Chen and Hill verbatim.; Verify 2606.31013 and 2609.09830 with metadata, and trace the contested f-002 and f-003 numbers to quoted excerpts: Eq. 13 (356 MeV), the f\_q formula, Eq. 38, the Table VII NLO rows, and the tree-level delta\_m0 formula (Eq. 9 / Ref. [28])..

\textbf{Wave 2} --- now exploring, in parallel: Pin down the pure-higgsino thermal mass from at least two independent published calculations, including recent work with bound-state formation, and quote the Planck 2018 Omega\_c h\textasciicircum{}2 verbatim. Confirm which mass the 1411.6930 pure-higgsino values (0.126/0.127) refer to.; Replace the linear 40 GeV extrapolation with direct reads of the LZ 2024 and XENONnT SI limits at about 1 TeV, a numeric xenon neutrino-fog value at about 1 TeV, and one next-generation projection (XLZD/DARWIN). Correct the f-008 arithmetic and quote Chen and Hill verbatim.; Verify 2606.31013 and 2609.09830 with metadata, and trace the contested f-002 and f-003 numbers to quoted excerpts: Eq. 13 (356 MeV), the f\_q formula, Eq. 38, the Table VII NLO rows, and the tree-level delta\_m0 formula (Eq. 9 / Ref. [28])..

\textbf{Finding:} Chen and Hill (1912.07795), verbatim from the Summary: ''The upper limit for the pure doublet is less than 10-48 cm2 in the TeV mass regime, consistent with the leading order estimate [18, 19], and lower than the estimated Xenon neutrino ?oor [26].'' Verbatim from Sec. 4: ''In the pure Higgsino limit, the cross section upper limit remains at or below 10 -48 cm2 for masses above a few hundred GeV.'' Abstract: ''The direct detection signal in the pure Higgsino limit remains below neutrino backgrounds for WIMPs in the TeV mass range.'' Sec. 4 also says that away from the pure limit, with c\_H/g2\textasciicircum{}2 of about 0.1, the TeV cross section is within reach of next-generation detectors, and with c\_H/g2\textasciicircum{}2 of about 0.01 it is close to the neutrino floor.

\textbf{Finding:} Correction to f-008: an upper bound below 1e-48 cm\textasciicircum{}2 compared with the LZ minimum of 2.2e-48 cm\textasciicircum{}2 (at 40 GeV) is a factor of more than 2.2 (more than 0.34 decades), not 2 orders of magnitude. The comparison that actually matters is at the same mass. At 1.1 TeV the linearly rescaled LZ limit of about 6e-47 cm\textasciicircum{}2, divided by the bound of 1e-48 cm\textasciicircum{}2, gives a ratio of at least about 60 (at least about 1.8 decades). That ratio is still an estimate: the LZ and XENONnT curves at 1 TeV were not read.

\textbf{Finding:} Planck 2018 (arXiv:1807.06209, abstract, page 1) states verbatim: ''A combined analysis gives dark matter density $\Omega$ch2 = 0.120$\pm$ 0.001'' (68\% confidence region), quoted for base-LambdaCDM from polarization, temperature and lensing combined.

\textbf{Finding:} Bottaro et al., ''The last Complex WIMPs standing'' (arXiv:2205.04486, 2022), is the independent recent calculation of thermal masses for complex electroweak n-plets, which includes the Y=1/2 doublet (the higgsino). Its companion paper arXiv:2107.09688 states that the method includes ''Sommerfeld enhancement and bound states formation at leading order'', and 2205.04486 says the doublet with hypercharge 1/2 is the notable exception that future direct detection cannot fully probe. The numeric doublet thermal mass was not read in this direction.

\textbf{Blocked --- Replace the linear 40 GeV extrapolation with direct reads of the LZ 2024 and XENONnT SI limits at about 1 TeV, a numeric xenon neutrino-fog value at about 1 TeV, and one next-generation projection (XLZD/DARWIN). Correct the f-008 arithmetic and quote Chen and Hill verbatim..} (The figure-extraction tool returned none of the SI limit plots from 2410.17036 or 2303.14729, and the tool budget ran out before the rest of the LZ text, the XENONnT text, HEPData or the XLZD Design Book (2410.17137) could be read.) Two things were established: the Chen and Hill quotes and the corrected f-008 arithmetic. The 1 TeV limits, the neutrino-fog value and the XLZD projection were not obtained. (1) Chen and Hill (arXiv:1912.07795), verbatim from the Summary: 'The upper limit for the pure doublet is less than 10\textasciicircum{}-48 cm2 in the TeV mass regime, consistent with the leading order estimate [18, 19], and lower than the estimated Xenon neutrino floor [26].' From Sec. 4: 'In the pure Higgsino limit, the cross section upper limit remains at or below 10\textasciicircum{}-48 cm2 for masses above a few hundred GeV.' Sec. 4 also says that with gaugino admixture c\_H/g2\textasciicircum{}2 of about 0.1 the TeV cross section is within reach of next-generation detectors, and with about 0.01 it is close to the neutrino floor. Finding f-010. These quotes are in the board finding only; the cite tool was not called in this direction, so they were not added to references.bib. (2) Correction to f-008: an upper bound below 1e-48 cm\textasciicircum{}2 compared with the LZ minimum of 2.2e-48 cm\textasciicircum{}2 at 40 GeV is a factor of more than 2.2 (more than 0.34 decades), not 2 orders of magnitude. Comparing at the same mass, the linearly rescaled LZ limit at 1.1 TeV (about 6e-47 cm\textasciicircum{}2, an estimate) is at least about 60 times the bound (at least about 1.8 decades). Finding f-011. (3) Not obtained: get\_paper\_figures did not return the limit plots (dead end x-01). For 2410.17036 it returned only Fig. 1 and Fig. S4, registered as Figures 10 and 11. For 2303.14729 it returned only calibration and spatial-distribution plots, registered as Figures 12-14. The LZ text chunk at characters 20000-28000 contained no 1 TeV value. The XLZD Design Book (arXiv:2410.17137) was identified but not read, so there are no fog or projection numbers yet (question q-07). The LZ/prediction, XENONnT/prediction and fog/prediction ratios at 1.1 TeV therefore remain estimates or are missing. Files in /Users/nramachandra/hep-genesis-testing/projects/ramp-research-higgsino-dd/runs/261016\_110437/: figures\_2410.17036/, figures\_2303.14729/, 1912.07795.txt, 2410.17036.txt.

\textbf{Blocked --- Pin down the pure-higgsino thermal mass from at least two independent published calculations, including recent work with bound-state formation, and quote the Planck 2018 Omega\_c h\textasciicircum{}2 verbatim. Confirm which mass the 1411.6930 pure-higgsino values (0.126/0.127) refer to..} (The tool budget ran out before the 2205.04486 full text (thermal-mass table) and 1411.6930 Sec. 3 could be read. A follow-up should download 2205.04486 and read its doublet thermal mass, and search 1411.6930.txt for the pure-higgsino mass statement.) One of the three success criteria is met. (1) Met: the Planck 2018 value is quoted verbatim from the abstract on page 1 of arXiv:1807.06209: 'A combined analysis gives dark matter density $\Omega$ch2 = 0.120$\pm$ 0.001' (68\% CL). It is at roughly characters 6500-7000 of /Users/nramachandra/hep-genesis-testing/projects/ramp-research-higgsino-dd/runs/261016\_110437/1807.06209.txt (finding f-012, tc-261016\_1-w2-d-04-82). It has not yet been added to references.bib with cite(). (2) Not met: a second independent calculation was identified but its number was not read. Bottaro et al., arXiv:2205.04486 ('The last Complex WIMPs standing', 2022), computes thermal masses for complex electroweak n-plets, including the Y=1/2 doublet (the higgsino). Its companion paper, arXiv:2107.09688, describes the method as including Sommerfeld enhancement and bound-state formation at leading order (finding f-013, medium confidence, abstracts only). Its doublet thermal mass was not read, so the spread against the 1411.6930-derived value of about 1.14 TeV (f-009) cannot be stated. For now, only one thermal-mass determination is available, and that one is a rescaled value. (3) Not met: the 1411.6930 mass attribution was not re-read in this direction. The 1172.31 GeV mass behind the pure-higgsino values 0.127/0.126 is still inferred, not stated verbatim (f-007 stays contested). The f-009 rescaling to 0.120 still gives about 1144 GeV (Sommerfeld) and 1139 GeV (perturbative). These numbers are derived, not published, and they depend on the inferred mass. One dead end was recorded: the abstract search for recent thermal-mass papers returned only unrelated LEP notes (x-02). Open question q-08 asks for the 2205.04486 doublet mass and its uncertainty.

\textbf{Finding:} Metadata confirms both IDs: arXiv:2606.31013 is 'Complementary Probes of Light Higgsinos: Electroweak Precision Measurements and Dark Matter Direct Detection' by Hamaguchi, Nagata and Osaki (2026), and arXiv:2609.09830 is 'Radiative Corrections to the Direct Detection of Inelastic Scattering of Higgsino-like Neutralino Dark Matter' by Chatterjee, Das, Pasha, Pukhov and Puri (2026). Both abstracts match how the survey uses them. This resolves q-02.

\textbf{Finding:} The full text of 2606.31013 supports the f-003 numbers. Eq. 13 reads Delta m+\textbar{}rad \textasciitilde{} (1/2) alpha2 MZ sin\textasciicircum{}2 thetaW (1 - 3MZ/(2 pi M\_chi+)), with the statement 'For M\_chi+ \textgreater{}\textgreater{} MZ, this mass splitting asymptotically approaches Delta m+\textbar{}rad \textasciitilde{} 356 MeV'. Eq. 33 reads f\_q \textasciitilde{} -(g\textasciicircum{}2/8M\_h\textasciicircum{}2)[1/(M2-\textbar{}mu\textbar{}) + tan\textasciicircum{}2 thetaW/(M1-\textbar{}mu\textbar{})](1 + sgn(mu) sin 2beta) and is valid for MZ \textless{}\textless{} M1,2-\textbar{}mu\textbar{}; the exact form is Eq. 32. The paper says the coupling is suppressed at large gaugino mass and 'for small tan beta and negative mu, due to an accidental cancellation'. Eq. 38 is sigma\_SI/SD,eff = (Omega\_H h\textasciicircum{}2/0.12) sigma\_SI/SD, computed with micrOMEGAs v6.3.0. The scan (Eq. 37) covers 150 GeV\textless{}\textbar{}mu\textbar{}\textless{}1000 GeV, 1.6\textless{}tan beta\textless{}50 and 0.5 TeV\textless{}M1\textless{}10 TeV, with M2=2M1 and M1\textgreater{}\textbar{}mu\textbar{}. Note: the f-003 label 'Eq. 13' is correct; the f\_q formula is Eq. 33, not unnumbered.

\textbf{Finding:} The full text of 2609.09830 supports the f-002 numbers. Table III gives one-loop on-shell delta = 36.37, 8.90, 16.39, 10.81 and 10.59 keV for BP1-BP5, with mu = -739.63, -606.64, -814.26, 1282.11 and 1400.0 GeV, M1 = -7830, -6292, -5838, 11354 and 12875 GeV, M2 = 17055, 42000, 30925, 12627 and 12050 GeV, and tan beta = 10.88-26.46. Table VII, NLO rows for BP1-BP5: elastic events 1.0280, 0.07729, 0.1176, 0.4127, 0.5074; inelastic events 0.00022, 0.07039, 0.0518, 0.1806, 0.1661; combined 1.02822, 0.14768, 0.1694, 0.5933, 0.6735; required exposure 9.41, 65.50, 57.10, 16.30, 14.36 t-y (LZ exposure 4.2 t-y; null result taken as N\_E \textless{} ln 10, about 2.3). Table VI for BP2 gives sigma\_inel NLO = 0.507e-12 pb against sigma\_el LO = 0.284e-12 pb (ratio 1.785). The paper describes this as 'approximately a factor of two larger than the elastic LO one'. The text says all BPs are consistent with LZ. Note: the BP1 scan sign pattern has M1\textless{}0 and M2\textgreater{}0 for mu=$\pm$0.5 or $\pm$1 TeV, with -21 TeV \textless{}= M1 \textless{}= -1 TeV and 1 TeV \textless{}= M2 \textless{}= 21 TeV (Eq. 25), yet BP4 and BP5 have M1 and M2 both positive.

\textbf{Finding:} 2609.09830 does not write out the tree-level delta\_m0 formula. It says only: 'as can be seen, for example, from Eq. (9) of Ref. [28], the corrections proportional to M\_Z\textasciicircum{}2/M1 and M\_Z\textasciicircum{}2/M2 depend on the relative signs of M1 and M2. In the parameter regions relevant for our analysis, these contributions can have opposite signs and partially cancel.' 2606.31013 also says only that the splittings are 'suppressed by M\_Z\textasciicircum{}2/M\_1,2'. The explicit form delta\_m0 \textasciitilde{} m\_Z\textasciicircum{}2 s\_W\textasciicircum{}2/M1 + m\_Z\textasciicircum{}2 c\_W\textasciicircum{}2/M2 therefore remains unsourced (q-05 stays open), and so does the estimate that same-sign M2=2M1 needs M1 of about 1.4e7 GeV. Counter-evidence to the 'tuned corners only' conclusion: BP4 and BP5 reach delta of about 10.6-10.8 keV with M1 and M2 both positive (about 11-13 TeV) and mu \textgreater{} 0. The O(10) keV splittings there come from on-shell one-loop corrections combined with relative signs between mu and the gauginos, not only from opposite-sign M1 and M2.

\textbf{Completed --- Verify 2606.31013 and 2609.09830 with metadata, and trace the contested f-002 and f-003 numbers to quoted excerpts: Eq. 13 (356 MeV), the f\_q formula, Eq. 38, the Table VII NLO rows, and the tree-level delta\_m0 formula (Eq. 9 / Ref. [28])..} Both papers are real, and every contested f-002/f-003 number except the tree-level delta\_m0 formula is now traced to a verbatim excerpt. That formula is still unsourced, so the M1 needed for delta \textless{} 364 keV was not re-derived. No cite() calls were made: the quotes are recorded on the board as f-015 and f-016, not in references.bib.

(1) IDs confirmed by metadata (f-014, closes q-02). 2606.31013 is Hamaguchi, Nagata and Osaki (2026), 'Complementary Probes of Light Higgsinos: Electroweak Precision Measurements and Dark Matter Direct Detection'. 2609.09830 is Chatterjee, Das, Pasha, Pukhov and Puri (2026), 'Radiative Corrections to the Direct Detection of Inelastic Scattering of Higgsino-like Neutralino Dark Matter'.

(2) 2606.31013 (f-015):

\begin{itemize}
  \item Eq. 13: Delta m+ \textasciitilde{} (1/2) alpha2 MZ sin\textasciicircum{}2 thetaW (1 - 3MZ/(2 pi M)), which tends to about 356 MeV for M \textgreater{}\textgreater{} MZ.
  \item The approximate f\_q formula is Eq. 33 (the exact form is Eq. 32). It is suppressed at large gaugino mass, and at small tan beta with mu \textless{} 0 through an accidental cancellation.
  \item Eq. 38: sigma\_eff = (Omega\_H h\textasciicircum{}2/0.12) sigma, computed with micrOMEGAs v6.3.0.
  \item Eq. 37 scan: 150 GeV \textless{} \textbar{}mu\textbar{} \textless{} 1000 GeV, 1.6 \textless{} tan beta \textless{} 50, 0.5 TeV \textless{} M1 \textless{} 10 TeV, with M2 = 2 M1.
\end{itemize}
(3) 2609.09830 (f-016):

\begin{itemize}
  \item Table III: one-loop on-shell delta = 36.37, 8.90, 16.39, 10.81 and 10.59 keV for BP1-BP5.
  \item Table VII, NLO events for BP1-BP5:
  \item inelastic: 0.00022, 0.07039, 0.0518, 0.1806, 0.1661
  \item elastic: 1.0280, 0.07729, 0.1176, 0.4127, 0.5074
  \item combined: 1.02822, 0.14768, 0.1694, 0.5933, 0.6735
  \item exposure needed to probe each point: 9.41, 65.50, 57.10, 16.30, 14.36 t-y (LZ has 4.2 t-y; the null-result limit used is N \textless{} ln 10, about 2.3).
  \item Table VI, BP2: NLO inelastic cross section 0.507e-12 pb against LO elastic 0.284e-12 pb, a ratio of 1.785.
\end{itemize}
(4) delta\_m0 formula (f-017). 2609.09830 only points to Eq. (9) of its Ref. [28]. It says the M\_Z\textasciicircum{}2/M1 and M\_Z\textasciicircum{}2/M2 terms depend on the relative signs of M1 and M2 and can partially cancel, but it does not write out the formula. So q-05 stays open, and so does the earlier estimate that same-sign M2 = 2 M1 needs M1 of about 1.4e7 GeV for delta \textless{} 364 keV.

(5) Correction to f-002 and the 'tuned corners only' conclusion: BP4 and BP5 reach about 10.6-10.8 keV with M1 and M2 both positive (about 11-13 TeV) and mu \textgreater{} 0. Opposite-sign gauginos are therefore not the only route to O(10) keV splittings; on-shell one-loop corrections and relative signs between mu and the gauginos also contribute.

Open: q-09, identify Ref. [28] of 2609.09830 and quote its Eq. (9).

Full texts: /Users/nramachandra/hep-genesis-testing/projects/ramp-research-higgsino-dd/runs/261016\_110437/2606.31013.txt and 2609.09830.txt.

\textbf{Critic, wave 2} --- continue: 7 findings stand up to the evidence, 1 questioned ('Chen and Hill (1912.07795), verbatim from the Summary: ''The upper limit for the pure doubl': Neither excerpt contains the quoted text. tc-...-76 (chars 0-3000) stops partway through the abstract, before the quoted abstract sentence. tc-...-81 (chars 22000-29000) covers the uncertainty methodology only. The Summary sentence ('less than 10\textasciicircum{}-48 cm2 ... lower than the estimated Xenon neutrino floor') and the Sec. 4 sentences, including the c\_H/g2\textasciicircum{}2 = 0.1 and 0.01 statements, are not shown. In a 37078-character document the Summary likely sits after char 29000, so it is probably outside both retrieved ranges. The quotes are plausibly correct and were earlier cited in d-02 (tc-...-64), but that call is not attached here. They need to be re-read with an excerpt that actually shows the text.).

Issues: The arithmetic is correct: 2.2e-48 / 1e-48 = 2.2, which is 0.34 decades, and 2.2e-48 x 1100/40 = 6.05e-47, giving a ratio of at least about 60 (about 1.78 decades). The key number is still a linear extrapolation from the 40 GeV minimum, not a read value. The 2.2e-48 input comes from f-005 and is not in this finding's excerpts. The headline 'where does the higgsino sit' ratio therefore still lacks a measured limit at 1 TeV.; The excerpt from tc-...-82 is truncated before the quoted sentence. The retrieved range (chars 5500-10000) does cover the stated location (about 6500-7000), and the value is standard, so it is accepted. It is still not in references.bib; cite() should be called.; None of the Table III, VI or VII numbers appear in the displayed excerpts. Only surrounding prose is shown, although the retrieved ranges (chars 27000-50000) plausibly contain the tables. This was the original reason f-002 was contested, so the trace is only partly verbatim. f-016 is accepted on the strength of the range coverage and the consistent BP3-BP5 prose. The numbers should still be recorded with cite() and table excerpts..

Stopping after wave 2: budget reached (2/2 waves used). Writing the report from the board.

\textit{Research report:}\par
\subsubsection*{The thermal pure higgsino (about 1.1 TeV) against Planck 2018, LZ 2024 and XENONnT: a literature survey}
\subsubsection*{Abstract}
The task was to compute the pure-higgsino relic density at several masses with inelastic-DM tools, and to place the 1.0--1.2 TeV thermal higgsino against LZ, XENONnT and the Planck abundance. No relic-density or inelastic-DM compute tool was available in this environment. No input data were available either. \textbf{Every Omega h\textasciicircum{}2 value below therefore comes from the literature, or is a labelled rescaling of a literature value; none was computed.}

The main findings are these:

\begin{itemize}
  \item \textbf{Relic density.} The one published pure-higgsino point found is Omega h\textasciicircum{}2 = 0.127 (perturbative) and 0.126 (Sommerfeld) in arXiv:1411.6930. The mass of about 1172 GeV attached to it is inferred from context, not stated. Rescaling that point as m\textasciicircum{}2 to the verbatim Planck 2018 value, Omega\_c h\textasciicircum{}2 = 0.120 $\pm$ 0.001, gives a thermal mass of about 1.14 TeV.
  \item \textbf{Elastic scattering.} The most complete elastic calculation (arXiv:1912.07795) bounds the pure-higgsino spin-independent (SI) cross section below 1e-48 cm\textasciicircum{}2 at TeV masses, below the xenon neutrino floor. LZ 2024 is tightest at 40 GeV, with 2.2e-48 cm\textasciicircum{}2. A linear extrapolation to 1.1 TeV, which is an estimate and not a read value, puts the LZ limit at about 6e-47 cm\textasciicircum{}2. That is at least about 60 times above the higgsino bound.
  \item \textbf{Inelastic scattering.} The kinematic threshold for up-scattering on xenon is about 364--368 keV at 1.0--1.1 TeV. Recent one-loop benchmarks with splittings of 9--36 keV predict at most about one LZ event and are all allowed.
\end{itemize}
Net conclusion: a thermal pure higgsino is not excluded by current direct detection. It becomes visible only through gaugino mixing or very small neutral splittings.

\subsubsection*{Background}
A pure higgsino is an SU(2) doublet with Y = 1/2. The literature places its thermal-relic mass at roughly 1.0--1.2 TeV, with the exact value set by the Sommerfeld and coannihilation treatment:

\begin{itemize}
  \item Beneke, Hellmann and Ruiz-Femenia, arXiv:1411.6930, give a Sommerfeld-enhanced coannihilation calculation.
  \item arXiv:2606.31013 quotes about 1 TeV.
  \item An earlier estimate, hep-ph/0610249, expected a Sommerfeld effect of about 10\% for higgsinos, against about 50\% for winos.
  \item A more recent independent framework includes bound-state formation as well as Sommerfeld enhancement (Bottaro et al., arXiv:2107.09688 and arXiv:2205.04486).
  \item The target abundance is from Planck 2018 (arXiv:1807.06209).
\end{itemize}
Pure-higgsino elastic scattering arises only at loop level and is suppressed by an accidental cancellation:

\begin{itemize}
  \item Hisano et al. give the leading-order result in arXiv:1104.0228, with gluon contributions in arXiv:1007.2601.
  \item Chen and Hill, arXiv:1912.07795, refine it with power corrections, QCD corrections to NNNLO and chiral-EFT nuclear effects.
  \item arXiv:2607.20588 (2026) argues that NLO corrections can push the cross section lower, or make it vanish, in parts of MSSM parameter space.
\end{itemize}
The experimental references are:

\begin{itemize}
  \item LZ, arXiv:2410.17036, with 4.2 tonne-years of exposure.
  \item XENONnT, arXiv:2303.14729.
  \item The ICRC 2025 review, arXiv:2510.17473, which notes that xenon experiments are entering the solar-neutrino-limited regime.
\end{itemize}
Two 2026 papers cover the visibility routes:

\begin{itemize}
  \item \textbf{Mixing (arXiv:2606.31013).} Hamaguchi, Nagata and Osaki show that gaugino mixing raises the Higgs coupling.
  \item \textbf{Inelastic scattering (arXiv:2609.09830).} Chatterjee et al. compute inelastic higgsino scattering at one loop for LZ.
\end{itemize}
Both IDs were confirmed by metadata lookup.

\subsubsection*{Results}
\paragraph{Planck target and the absence of a compute tool}
Planck 2018 states verbatim, on page 1 of the abstract: ''A combined analysis gives dark matter density $\Omega$ch2 = 0.120$\pm$ 0.001'' (68\% CL, base $\Lambda$CDM). This value was read from the downloaded text but not added to the bibliography file.

A file listing returned an empty input set. The tool catalog contains only statistics/histogram, test-job and literature tools. The task's request to ''compute the relic density'' therefore could not be carried out.

\paragraph{Relic density against mass}
arXiv:1411.6930 (Sec. 3) reports the following:

\begin{itemize}
  \item \textbf{Pure-higgsino model} (sfermions and heavy Higgses decoupled): Omega h\textasciicircum{}2 = 0.127 without and 0.126 with Sommerfeld enhancement, including off-diagonal rates. The net Sommerfeld effect is below 1\%, because enhancement in the neutral channels cancels suppression in the neutral--charged channels.
  \item \textbf{Higgsino-like pMSSM point 1627006} (higgsino fraction 0.98, $\Delta$m+ = 1.8 GeV, $\Delta$m0 = 9.5 GeV): Omega h\textasciicircum{}2 = 0.108 perturbative and 0.100 with Sommerfeld, an 8\% reduction.
  \item \textbf{Normalisation and accuracy.} The paper quotes about 5\% Sommerfeld accuracy at best and normalises to Planck 2013 (0.1187 $\pm$ 0.0017).
\end{itemize}
The pure-higgsino mass of 1172.31 GeV is the mass of the pMSSM comparison point. Its attribution to the pure-higgsino numbers was inferred from context and is not stated in the paper. The critic flagged this, and it remains unconfirmed.

[[FIGURE 6]]

Figure 6 reproduces Fig. 7 of arXiv:1411.6930. It shows the yield ratio Y/Y\_pert, Sommerfeld over perturbative, for the pMSSM higgsino-like point with and without off-diagonal rates, and for the pure-higgsino model. It shows the size of the Sommerfeld correction, not Omega h\textasciicircum{}2 against mass. No published Omega h\textasciicircum{}2(m) curve was found.

The table below is an m\textasciicircum{}2 rescaling of the single point (0.126 $\times$ (m/1172.31 GeV)\textasciicircum{}2). It is an interpolation, not published values, and it depends on the inferred mass.

\begin{longtable}{@{}p{0.160\textwidth}p{0.160\textwidth}p{0.160\textwidth}p{0.160\textwidth}p{0.160\textwidth}p{0.160\textwidth}@{}}
\toprule
\textbf{m (TeV)} & \textbf{0.9} & \textbf{1.0} & \textbf{1.1} & \textbf{1.2} & \textbf{1.3} \\ \midrule \endhead
Omega h\textasciicircum{}2 (rescaled) & 0.074 & 0.092 & 0.111 & 0.132 & 0.155 \\
\bottomrule\end{longtable}
Omega h\textasciicircum{}2 = 0.120 falls at about 1144 GeV (Sommerfeld) or 1139 GeV (perturbative), that is about 1.14 TeV.

A second, independent thermal mass was identified but not read: the doublet entry of arXiv:2205.04486. That paper also states that the Y = 1/2 doublet is the notable case that future direct detection cannot fully probe.

\paragraph{Elastic SI cross section against LZ and XENONnT}
\textbf{Prediction.} Hisano et al. (arXiv:1104.0228) find an SI cross section of O(1e-46--1e-48) cm\textasciicircum{}2 at leading order for pure electroweak multiplets, nearly mass-independent above about 1 TeV. The positive twist-2 terms cancel against the negative Higgs, gluon and scalar terms. The critic noted that the numeric range was not visible in the excerpts read. The calculation scans the Higgs mass and predates its measurement.

Chen and Hill (arXiv:1912.07795) state verbatim: ''The upper limit for the pure doublet is less than 10\textasciicircum{}-48 cm2 in the TeV mass regime, consistent with the leading order estimate [18, 19], and lower than the estimated Xenon neutrino floor [26].'' Sec. 4 adds that the bound ''remains at or below 10\textasciicircum{}-48 cm2 for masses above a few hundred GeV''. The same section says that with gaugino admixture c\_H/g2\textasciicircum{}2 of about 0.1 the TeV rate comes within next-generation reach, and with about 0.01 it lies near the floor. These quotes were recorded, but the critic did not treat them as fully verified, and only the earlier citation entry is in the bibliography. Only an upper bound was obtained, not a central value with an uncertainty.

\textbf{Limits.} LZ 2024 gives its strongest SI exclusion as 2.2e-48 cm\textasciicircum{}2 (90\% CL), with a best median sensitivity of 5.1e-48 cm\textasciicircum{}2, both at 40 GeV. No 1 TeV value was read.

\textbf{Comparison.} Comparing the 1e-48 cm\textasciicircum{}2 bound with the 40 GeV minimum gives a factor of only more than 2.2 (more than 0.34 decades). An earlier ''two orders of magnitude'' statement was wrong on this point and has been corrected.

At equal mass, scaling the LZ limit linearly gives 2.2e-48 $\times$ 1100/40 $\approx$ 6e-47 cm\textasciicircum{}2 at 1.1 TeV. That is a ratio of at least about 60 (at least about 1.8 decades) above the higgsino bound. This is an extrapolation. The XENONnT 1 TeV limit, the LZ 1 TeV curve value and a numeric neutrino-fog value were not obtained.

\paragraph{Inelastic threshold and LZ inelastic sensitivity}
\textbf{Threshold.} The threshold was worked out by hand:

\begin{itemize}
  \item m\_Xe = 131 $\times$ 0.9315 GeV = 122.0 GeV.
  \item $\mu$\_A = 108.7 GeV at 1.0 TeV and 109.8 GeV at 1.1 TeV.
  \item v\_max = 544 + 232 = 776 km/s, so $\beta$\textasciicircum{}2 = 6.70e-6.
  \item $\delta$\_max = $\mu$\_A $\beta$\textasciicircum{}2/2 gives 364 keV at 1.0 TeV and 368 keV at 1.1 TeV.
\end{itemize}
The relation v\_min = sqrt(2$\delta$/$\mu$\_A) is Eq. A7 of arXiv:2609.09830. The result depends on the escape speed and Earth velocity assumed.

\textbf{Benchmarks.} arXiv:2609.09830 does not give an LZ-excluded $\delta$\_m0 range. Its five benchmarks are:

\begin{itemize}
  \item Table III: $\mu$ = -739.63, -606.64, -814.26, 1282.11 and 1400.0 GeV; M1 = -7830, -6292, -5838, 11354 and 12875 GeV; M2 = 17055, 42000, 30925, 12627 and 12050 GeV; tan $\beta$ = 10.88--26.46.
  \item One-loop on-shell $\delta$ = 36.37, 8.90, 16.39, 10.81 and 10.59 keV for BP1--BP5.
\end{itemize}
Table VII gives the NLO LZ event counts:

\begin{longtable}{@{}p{0.160\textwidth}p{0.160\textwidth}p{0.160\textwidth}p{0.160\textwidth}p{0.160\textwidth}p{0.160\textwidth}@{}}
\toprule
\textbf{} & \textbf{BP1} & \textbf{BP2} & \textbf{BP3} & \textbf{BP4} & \textbf{BP5} \\ \midrule \endhead
Elastic & 1.0280 & 0.07729 & 0.1176 & 0.4127 & 0.5074 \\
Inelastic & 0.00022 & 0.07039 & 0.0518 & 0.1806 & 0.1661 \\
Combined & 1.02822 & 0.14768 & 0.1694 & 0.5933 & 0.6735 \\
Exposure needed (t-y) & 9.41 & 65.50 & 57.10 & 16.30 & 14.36 \\
\bottomrule\end{longtable}
For comparison, LZ has 4.2 t-y, and the null-result criterion used is N \textless{} ln 10 $\approx$ 2.3. All benchmarks are consistent with LZ.

NLO vertex corrections can exceed 100\%. For BP2, the NLO inelastic cross section is 0.507e-12 pb, against an LO elastic cross section of 0.284e-12 pb (ratio 1.785).

The critic noted that the table rows were not visible in the displayed excerpts, although the ranges read cover them. They are accepted on that basis.

\paragraph{Splittings and mixing}
\textbf{Charged--neutral splitting.} Eq. 13 of arXiv:2606.31013 gives the pure-limit radiative splitting as $\Delta$m+ $\approx$ (1/2) $\alpha$2 MZ sin\textasciicircum{}2$\theta$W (1 - 3MZ/(2$\pi$ M)). This tends to about 356 MeV for M ? MZ, rather than the 340 MeV in the brief.

\textbf{Higgs coupling.} Eq. 33 (approximate form of Eq. 32) gives f\_q $\approx$ -(g\textasciicircum{}2/8M\_h\textasciicircum{}2)[1/(M2-\textbar{}$\mu$\textbar{}) + tan\textasciicircum{}2$\theta$W/(M1-\textbar{}$\mu$\textbar{})](1 + sgn($\mu$) sin 2$\beta$). This is suppressed at large gaugino mass, and at small tan $\beta$ with $\mu$ \textless{} 0 through an accidental cancellation.

\textbf{Scan.} The scan in Eq. 37 covers 150 \textless{} \textbar{}$\mu$\textbar{} \textless{} 1000 GeV, 1.6 \textless{} tan $\beta$ \textless{} 50 and 0.5 \textless{} M1 \textless{} 10 TeV, with M2 = 2 M1. Rates are rescaled as $\sigma$\_eff = (Omega\_H h\textasciicircum{}2/0.12) $\sigma$ (Eq. 38, micrOMEGAs v6.3.0). arXiv:2609.09830 likewise weakens its limits proportionally for sub-thermal fractions.

Direct detection is sensitive mainly where mixing produces sizable charged--neutral splittings, and in that region it reaches the about 1 TeV thermal mass. The reach figures of that paper were not extracted.

\textbf{Neutral splitting $\delta$\_m0.} The tree-level formula was not sourced. arXiv:2609.09830 refers only to ''Eq. (9) of Ref. [28]'' and states that the M\_Z\textasciicircum{}2/M1 and M\_Z\textasciicircum{}2/M2 terms can partially cancel. An earlier estimate, that same-sign M2 = 2 M1 needs M1 $\approx$ 1.4e7 GeV for $\delta$ \textless{} 364 keV, rests on that unsourced formula and should be disregarded. The critic noted that a same-sign tree-level splitting at about 12 TeV would be only about 0.5--0.7 MeV.

BP4 and BP5 reach about 10.6--10.8 keV with M1 and M2 both positive (about 11--13 TeV) and $\mu$ \textgreater{} 0. So the earlier claim that only opposite-sign gauginos give O(10) keV splittings was revised. On-shell one-loop corrections and the relative signs between $\mu$ and the gauginos also contribute.

\subsubsection*{Discussion}
The survey supports the brief's qualitative picture.

\begin{itemize}
  \item \textbf{Thermal mass.} A thermal pure higgsino sits near 1.1 TeV. The single rescaled point gives about 1.14 TeV, inside the 1.0--1.2 TeV window. The Sommerfeld effect is negligible in the pure limit (below 1\%), smaller than the earlier 10\% expectation of hep-ph/0610249.
  \item \textbf{Elastic scattering.} The elastic SI rate is bounded below 1e-48 cm\textasciicircum{}2. By the extrapolation above, it lies roughly two decades below the current LZ reach at 1.1 TeV and below the xenon floor. Pure-higgsino elastic scattering is therefore not a viable discovery channel for current or, per arXiv:2205.04486 and arXiv:1912.07795, future xenon detectors.
  \item \textbf{Visibility routes.} The higgsino becomes visible through gaugino mixing, which raises f\_q and the splittings (arXiv:2606.31013). It can also become visible through neutral splittings well below the about 365 keV kinematic limit. Even at 9--36 keV, however, the one-loop benchmarks of arXiv:2609.09830 give fewer than about 1.03 LZ events and need 9.4--65.5 t-y.
\end{itemize}
Two points of tension remain:

\begin{itemize}
  \item \textbf{Splitting value.} The brief's 340 MeV charged--neutral splitting differs from the 356 MeV one-loop asymptote quoted here. The critic suggested that a two-loop value of about 341 MeV may reconcile them, but this was not read.
  \item \textbf{Robustness of the cancellation.} Whether the elastic cancellation is robust is unresolved. arXiv:2607.20588 suggests NLO effects can drive the rate lower still.
\end{itemize}
Dead ends:

\begin{itemize}
  \item \textbf{Limit figures.} Figure extraction from arXiv:2410.17036 and arXiv:2303.14729 returned no limit plots, only data-distribution and calibration figures. The LZ text region read had no 1 TeV value.
  \item \textbf{Abstract search.} An abstract search for recent thermal-mass papers returned unrelated LEP notes. Direct arXiv-ID lookup proved more effective.
\end{itemize}
\subsubsection*{Limitations}
\begin{itemize}
  \item \textbf{No computation.} No compute tool existed, so no relic density, event rate or reach map was computed. The m\textasciicircum{}2 table is an interpolation from one point with an inferred mass.
  \item \textbf{Thermal-mass spread.} Only one thermal-mass determination was available. The bound-state-inclusive value of arXiv:2205.04486 was not read, so the spread across calculations is unknown.
  \item \textbf{Direct-detection numbers not read.} The LZ and XENONnT limits at 1 TeV, the neutrino-fog value at 1 TeV, and next-generation projections (XLZD, arXiv:2410.17137) were not read. The headline ratio of about 60 or more remains an extrapolation.
  \item \textbf{$\delta$\_m0 formula.} The tree-level $\delta$\_m0 formula is unsourced.
  \item \textbf{Indirect detection.} Indirect detection was not addressed.
  \item \textbf{Bibliography.} Several quotes (Planck, the Chen and Hill passages, the Table III/VI/VII traces) are recorded on the board but not cited into references.bib.
  \item \textbf{Budget.} Both waves were used.
\end{itemize}
\subsubsection*{Open directions}
\begin{enumerate}
  \item Read the Y = 1/2 doublet thermal mass and its uncertainty from arXiv:2205.04486, to give a second, bound-state-inclusive determination to compare with about 1.14 TeV.
  \item Obtain the LZ 2024 and XENONnT SI limits at 1 TeV from HEPData or the remaining paper text, replacing the linear extrapolation.
  \item Read the XLZD Design Book (arXiv:2410.17137) for the xenon neutrino-fog value and the projected SI sensitivity at about 1 TeV.
  \item Confirm in the text of arXiv:1411.6930 which mass the pure-higgsino values 0.126/0.127 refer to.
  \item Identify Ref. [28] of arXiv:2609.09830 and quote its Eq. (9), to map $\delta$\_m0 against ($\mu$, M1, M2) with explicit sign dependence.
  \item Add one sourced sentence on indirect-detection (for example H.E.S.S. Galactic Center) limits on TeV higgsinos.
  \item Provide a relic-density compute tool (for example micrOMEGAs) so that Omega h\textasciicircum{}2(m) at 0.9--1.3 TeV can actually be computed.
\end{enumerate}
\end{document}
